\begin{abstract}
Hand-designed rewards and default settings shape what a reinforcement-learning agent learns, yet realistic problems offer no optimum to compare against. This paper studies a shelf-restocking robot in a store small enough to solve exactly, and uses the exact optimum to measure both the reward and the learners. The state is reduced to 768 Markov states, an exact model is validated against the simulator, and every policy, tabular or neural, is evaluated exactly on ten store layouts. Dynamic programming shows that the project's original reward pays for walking back and forth near a shelf. Its optimal policy travels 6\% further on average (up to 17\%), at a cost of only 0.3\% of the objective, whereas potential-based versions of the same guidance leave the optimal policy unchanged. Among nine learners trained under four rewards with identical budgets, DQN and Double DQN come closest to the optimum (normalised score 0.96), ahead of PPO (0.79) and every tabular method (at most 0.73). Both bonuses in the original reward help DQN, and their potential-based forms recover about half of that gain. For tabular learners the initial Q-value matters more than the reward: since potential-based shaping is equivalent to initialisation, moving the starting value from 0 to $-150$ triples Q-Learning's score, and direct counts reveal the failure modes on either side. Finally, REINFORCE's argmax policy scores less than half of its stochastic policy, so policy-gradient agents must be evaluated as stochastic policies.
\end{abstract}

\begin{IEEEkeywords}
reinforcement learning, reward shaping, dynamic programming, Q-value initialisation, deep Q-network, retail automation, inventory restocking
\end{IEEEkeywords}
